% ------------------------------------------------
% THỰC NGHIỆM (Thành viên A). Ghép vào báo cáo chính bằng % ------------------------------------------------
% THỰC NGHIỆM (Thành viên A). Ghép vào báo cáo chính bằng % ------------------------------------------------
% THỰC NGHIỆM (Thành viên A). Ghép vào báo cáo chính bằng % ------------------------------------------------
% THỰC NGHIỆM (Thành viên A). Ghép vào báo cáo chính bằng \input{sections/experiments}
% Cần: booktabs, multirow (\usepackage{multirow}), tham khảo trong custom.bib (zalo2020traffic, lin2014coco, jocher2023ultralytics,
% ren2015faster, lin2017fpn, akyon2022sahi, efron1993bootstrap)
% ------------------------------------------------

\section{Thực nghiệm}
\label{sec:experiments}

\subsection{Dữ liệu và cách chia tập}
\label{sec:data-split}

\paragraph{Gộp box trùng.}
Khi kiểm tra nhãn bằng mắt, nhóm phát hiện nhiều biển bị gán 2--3 box gần như trùng nhau, cùng lớp, chỉ lệch 1--2 pixel. Nếu giữ nguyên, mỗi biển chỉ khớp được với một dự đoán, nên các box lặp luôn bị tính là bỏ sót. Với tập test ban đầu, recall tối đa chỉ còn khoảng 77\% kể cả khi mô hình hoàn hảo.
Vì vậy, trong mỗi ảnh, nhóm gộp các box \emph{cùng lớp} có IoU $\geq 0{,}45$ thành một box có tọa độ trung bình. Các box được nối bắc cầu: A trùng B, B trùng C thì A, B, C chung một cụm.
Ngưỡng này an toàn vì hai biển thật đứng cạnh nhau gần như không chồng lên nhau (IoU $\approx 0$): 124/125 cặp có IoU trong khoảng $[0{,}45; 0{,}7)$ là box lặp cùng lớp.
Sau khi gộp, số box giảm từ 11.000 xuống 9.077 (17,5\%). Tập test bị nặng nhất: 24,7\% box bị bỏ.

\paragraph{Đặc điểm dữ liệu.}
Sau khi gộp, 77,6\% box là box nhỏ theo chuẩn COCO (diện tích $< 32^2$ pixel). Cạnh trung vị $\sqrt{wh}$ chỉ 16,5 pixel trên ảnh rộng 1622 pixel.
Khi thu nhỏ ảnh về cạnh dài 640, cạnh trung vị còn 6,5 pixel, và 68,9\% box có cạnh dưới 10 pixel.
Lớp nhiều nhất (\textit{Nguy hiểm}, 2.329 box) nhiều gấp 4,8 lần lớp ít nhất (\textit{Cấm rẽ}, 489 box).
Mỗi lớp đã gộp cả biến thể trái và phải, nên lật ngang ảnh không làm đổi lớp.

\paragraph{Chia tập theo tuyến đường.}
Ảnh là các khung hình liên tiếp trên 38 tuyến đường (trường \texttt{street\_id}). Nếu chia ngẫu nhiên theo ảnh, hai khung hình gần như giống hệt nhau có thể rơi vào cả train lẫn tập đánh giá.
Vì vậy nhóm chia theo tuyến đường, tỉ lệ 70/15/15 (seed 42): mọi ảnh của một tuyến đường nằm trọn trong một tập. Các tuyến đường được phân tầng theo lớp hiếm nhất có trong tuyến, để lớp hiếm có mặt ở cả ba tập (Bảng~\ref{tab:split}).
Mục~
ef{sec:leakage} định lượng mức đánh giá quá cao nếu chia ngẫu nhiên.
Mọi mô hình dùng chung bộ chia này. Chọn cấu hình chỉ dựa trên tập val. Tập test chỉ được chấm \emph{một lần}, cho các mô hình cuối cùng.

\begin{table}[t]
\centering
\small
\begin{tabular}{lrrr}
\toprule
Tập & Ảnh & Box & Tuyến đường \\
\midrule
Train & 2.933 & 5.693 & 28 \\
Val   &   696 & 1.636 &  5 \\
Test  &   871 & 1.748 &  5 \\
\bottomrule
\end{tabular}
\caption{Ba tập dữ liệu sau khi gộp box trùng, chia theo tuyến đường.}
\label{tab:split}
\end{table}

\subsection{Thiết lập huấn luyện}
\label{sec:setup}

Cả hai họ mô hình đều khởi tạo từ trọng số tiền huấn luyện trên COCO \citep{lin2014coco} và được huấn luyện trên GPU Tesla T4 (Kaggle).

Các cấu hình được mô tả ở Mục~\ref{sec:method}. Phần này chỉ liệt kê siêu tham số.

\paragraph{YOLOv8s} (Ultralytics 8.4.204).
\begin{itemize}
\item 50 epoch, dừng sớm nếu 15 epoch liên tiếp không cải thiện, seed 0.
\item Batch 16 ở imgsz 640, batch 8 ở imgsz 1280.
\item Augmentation để mặc định: mosaic (tắt ở 10 epoch cuối), lật ngang với xác suất 0,5, co giãn $\pm 50\%$.
\item Cấu hình cắt ô: 40 epoch, batch 16. Một box chỉ được giữ trong ô nếu ít nhất 60\% diện tích nằm trong ô. Các ô không có biển chỉ giữ ngẫu nhiên 15\%.
\item Thời gian: mỗi lần huấn luyện ở imgsz 1280 mất khoảng 2,2 giờ trên T4.
\end{itemize}

\paragraph{Faster R-CNN} (torchvision).
\begin{itemize}
\item SGD: lr 0,0025, momentum 0,9, weight decay $5 \cdot 10^{-4}$.
\item 30 epoch, giảm lr 10 lần ở epoch 20 và 26.
\item Batch 8, FP16, cắt gradient ở chuẩn 1,0.
\end{itemize}
Thời gian huấn luyện của các cấu hình còn lại không được ghi lại.

\subsection{Chỉ số đánh giá}
\label{sec:metrics}

Mọi mô hình xuất dự đoán theo cùng một định dạng: tọa độ pixel trên ảnh gốc, ngưỡng độ tin cậy rất thấp (0,001), để không cắt mất phần cuối của đường precision--recall.
Tất cả được chấm bằng cùng một script dựa trên \texttt{pycocotools}.
Chỉ số chính là mAP@0.5. Nhóm báo cáo thêm mAP@0.5:0.95 và AP theo kích thước theo chuẩn COCO: nhỏ ($< 32^2$), vừa ($32^2$--$96^2$) và lớn ($> 96^2$ pixel).
Tập val chỉ có 23 biển lớn, nên AP\textsubscript{L} trên val dao động mạnh và không được dùng để so sánh.
Để biết chênh lệch giữa hai mô hình có ý nghĩa thống kê hay không, nhóm dùng \emph{paired bootstrap} \citep{efron1993bootstrap}: lấy mẫu lại có hoàn lại các ảnh của tập đánh giá 1000 lần, mỗi lần tính lại AP50 của các mô hình trên cùng bộ ảnh, rồi lấy khoảng 2,5\%--97,5\% của chênh lệch làm khoảng tin cậy 95\%.

\subsection{Kết quả}
\label{sec:results}

\begin{table*}[t]
\centering
\small
\begin{tabular}{llcccccccc}
\toprule
& & \multicolumn{4}{c}{Val (696 ảnh)} & \multicolumn{4}{c}{Test (871 ảnh)} \\
\cmidrule(lr){3-6} \cmidrule(lr){7-10}
Mô hình & Cấu hình & mAP@.5 & mAP@.5:.95 & AP\textsubscript{S} & AP\textsubscript{M} & mAP@.5 & mAP@.5:.95 & AP\textsubscript{S} & AP\textsubscript{M} \\
\midrule
\multirow{4}{*}{YOLOv8s}
& 640 (baseline)    & 61,7 & 39,8 & 29,8 & 70,5 & -- & -- & -- & -- \\
& 640 + SAHI        & 78,1 & 45,2 & 37,3 & 70,8 & -- & -- & -- & -- \\
& 1280              & 82,2 & 56,6 & 50,2 & \textbf{76,1} & 76,8 & 50,8 & 45,1 & \textbf{76,9} \\
& Cắt ô 640 + SAHI  & \textbf{83,9} & \textbf{57,9} & \textbf{52,3} & 75,2 & \textbf{78,1} & \textbf{51,4} & \textbf{46,8} & 75,1 \\
\midrule
\multirow{3}{*}{\shortstack[l]{Faster R-CNN\\R50-FPN}}
& exp01 (baseline)   & 64,8 & 41,6 & 33,5 & 68,1 & -- & -- & -- & -- \\
& exp02 (anchor nhỏ) & 72,5 & 47,7 & 40,3 & 71,2 & -- & -- & -- & -- \\
& exp03 (+ augment)  & 75,2 & 48,6 & 41,3 & 71,8 & 69,3 & 41,8 & 35,9 & 71,4 \\
\bottomrule
\end{tabular}
\caption{Kết quả trên val và test (\%). Chỉ các mô hình cuối cùng được chấm trên test, và chỉ chấm một lần.
Dự đoán val của Faster R-CNN được xuất ở ngưỡng 0,05 thay vì 0,001, nên mAP trên val của Faster R-CNN có thể thấp hơn thực tế một chút; dự đoán test được xuất ở ngưỡng 0,0001.}
\label{tab:results}
\end{table*}

Bảng~\ref{tab:results} tổng hợp kết quả. Ba nhận xét chính:

\paragraph{Biển nhỏ quyết định kết quả.}
Ở cấu hình baseline, AP của biển nhỏ kém biển vừa 40,7 điểm với YOLO (ảnh 640) và 34,6 điểm với Faster R-CNN.
Các cách giữ được độ phân giải chủ yếu cải thiện biển nhỏ. Với YOLO, chuyển từ 640 lên 1280 tăng AP\textsubscript{S} từ 29,8 lên 50,2 (+20,4), trong khi AP\textsubscript{M} chỉ tăng 5,6.
Với Faster R-CNN, anchor nhỏ và không thu nhỏ ảnh (exp02) tăng AP\textsubscript{S} thêm 6,8 điểm.
SAHI áp lên mô hình 640 có sẵn (không huấn luyện lại) đã tăng mAP@0.5 thêm 16,4 điểm. Nhưng mAP@0.5:0.95 chỉ tăng 5,4 điểm, nghĩa là box ghép từ các ô kém khít hơn.

\paragraph{YOLOv8s tốt hơn Faster R-CNN rõ rệt.}
Trên test, YOLOv8s 1280 hơn Faster R-CNN exp03 7,5 điểm mAP@0.5. Khoảng tin cậy 95\% của chênh lệch là $[4{,}5; 10{,}2]$, tức có ý nghĩa thống kê.
Faster R-CNN kém hơn ở cả 7 lớp, trong đó 4 lớp có chênh lệch có ý nghĩa thống kê. Kém rõ nhất là \textit{Cấm còn lại} ($-20{,}2$) và \textit{Cấm rẽ} ($-10{,}2$).

\paragraph{Cắt ô + SAHI và 1280 ngang nhau.}
Trên test, cắt ô + SAHI hơn 1280 1,3 điểm mAP@0.5. Nhưng khoảng tin cậy $[-0{,}5; +3{,}2]$ chứa 0, nên chưa đủ bằng chứng để nói cách nào tốt hơn.
Cắt ô + SAHI tốt hơn ở \textit{Hiệu lệnh} (+7,7) và \textit{Cấm còn lại} (+5,1), nhưng tạo nhiều box lặp hơn (78 so với 17) và báo nhầm nền nhiều hơn (123 so với 97 box).
Vì vậy cả hai cấu hình đều được giữ làm kết quả cuối. Cấu hình 1280 đơn giản hơn và suy luận nhanh hơn.

\paragraph{Từ val sang test.}
mAP@0.5 của cả ba mô hình cuối đều giảm 5--6 điểm từ val sang test, và thứ hạng giữa chúng không đổi.
Tập test khó hơn val: 21\% biển trong test nhỏ hơn 8 pixel, trong khi ở val chỉ có 10\%. Ở cả ba mô hình, nhóm biển dưới 8 pixel bị bỏ sót 60--78\%.
Chọn cấu hình trên val vẫn dẫn tới cùng kết luận trên test.

\begin{table}[t]
\centering
\small
\begin{tabular}{lrccc}
\toprule
Lớp & Biển & 1280 & Cắt ô+SAHI & FRCNN \\
\midrule
Cấm ngược chiều & 253 & 84,7 & \textbf{86,8} & 84,2 \\
Cấm dừng và đỗ  & 209 & 87,9 & \textbf{90,1} & 82,7 \\
Cấm rẽ          &  81 & \textbf{64,6} & 58,9 & 54,4 \\
Giới hạn tốc độ &  79 & \textbf{65,4} & 61,7 & 58,4 \\
Cấm còn lại     & 176 & 70,0 & \textbf{75,1} & 49,9 \\
Nguy hiểm       & 690 & 87,6 & \textbf{89,1} & 81,1 \\
Hiệu lệnh       & 260 & 77,5 & \textbf{85,3} & 74,7 \\
\bottomrule
\end{tabular}
\caption{AP50 (\%) theo lớp trên test. FRCNN là Faster R-CNN exp03.}
\label{tab:per-class-test}
\end{table}

Theo Bảng~\ref{tab:per-class-test}, ba lớp biển tròn viền đỏ (\textit{Cấm rẽ}, \textit{Giới hạn tốc độ}, \textit{Cấm còn lại}) có AP thấp nhất ở cả ba mô hình. \textit{Cấm ngược chiều} thì ngược lại: lớp này có tới 90,5\% box nhỏ nhưng vẫn đạt AP cao. Mục~
ef{sec:confusion} giải thích hai hiện tượng này.

% Cần: booktabs, multirow (\usepackage{multirow}), tham khảo trong custom.bib (zalo2020traffic, lin2014coco, jocher2023ultralytics,
% ren2015faster, lin2017fpn, akyon2022sahi, efron1993bootstrap)
% ------------------------------------------------

\section{Thực nghiệm}
\label{sec:experiments}

\subsection{Dữ liệu và cách chia tập}
\label{sec:data-split}

\paragraph{Gộp box trùng.}
Khi kiểm tra nhãn bằng mắt, nhóm phát hiện nhiều biển bị gán 2--3 box gần như trùng nhau, cùng lớp, chỉ lệch 1--2 pixel. Nếu giữ nguyên, mỗi biển chỉ khớp được với một dự đoán, nên các box lặp luôn bị tính là bỏ sót. Với tập test ban đầu, recall tối đa chỉ còn khoảng 77\% kể cả khi mô hình hoàn hảo.
Vì vậy, trong mỗi ảnh, nhóm gộp các box \emph{cùng lớp} có IoU $\geq 0{,}45$ thành một box có tọa độ trung bình. Các box được nối bắc cầu: A trùng B, B trùng C thì A, B, C chung một cụm.
Ngưỡng này an toàn vì hai biển thật đứng cạnh nhau gần như không chồng lên nhau (IoU $\approx 0$): 124/125 cặp có IoU trong khoảng $[0{,}45; 0{,}7)$ là box lặp cùng lớp.
Sau khi gộp, số box giảm từ 11.000 xuống 9.077 (17,5\%). Tập test bị nặng nhất: 24,7\% box bị bỏ.

\paragraph{Đặc điểm dữ liệu.}
Sau khi gộp, 77,6\% box là box nhỏ theo chuẩn COCO (diện tích $< 32^2$ pixel). Cạnh trung vị $\sqrt{wh}$ chỉ 16,5 pixel trên ảnh rộng 1622 pixel.
Khi thu nhỏ ảnh về cạnh dài 640, cạnh trung vị còn 6,5 pixel, và 68,9\% box có cạnh dưới 10 pixel.
Lớp nhiều nhất (\textit{Nguy hiểm}, 2.329 box) nhiều gấp 4,8 lần lớp ít nhất (\textit{Cấm rẽ}, 489 box).
Mỗi lớp đã gộp cả biến thể trái và phải, nên lật ngang ảnh không làm đổi lớp.

\paragraph{Chia tập theo tuyến đường.}
Ảnh là các khung hình liên tiếp trên 38 tuyến đường (trường \texttt{street\_id}). Nếu chia ngẫu nhiên theo ảnh, hai khung hình gần như giống hệt nhau có thể rơi vào cả train lẫn tập đánh giá.
Vì vậy nhóm chia theo tuyến đường, tỉ lệ 70/15/15 (seed 42): mọi ảnh của một tuyến đường nằm trọn trong một tập. Các tuyến đường được phân tầng theo lớp hiếm nhất có trong tuyến, để lớp hiếm có mặt ở cả ba tập (Bảng~\ref{tab:split}).
Mục~
ef{sec:leakage} định lượng mức đánh giá quá cao nếu chia ngẫu nhiên.
Mọi mô hình dùng chung bộ chia này. Chọn cấu hình chỉ dựa trên tập val. Tập test chỉ được chấm \emph{một lần}, cho các mô hình cuối cùng.

\begin{table}[t]
\centering
\small
\begin{tabular}{lrrr}
\toprule
Tập & Ảnh & Box & Tuyến đường \\
\midrule
Train & 2.933 & 5.693 & 28 \\
Val   &   696 & 1.636 &  5 \\
Test  &   871 & 1.748 &  5 \\
\bottomrule
\end{tabular}
\caption{Ba tập dữ liệu sau khi gộp box trùng, chia theo tuyến đường.}
\label{tab:split}
\end{table}

\subsection{Thiết lập huấn luyện}
\label{sec:setup}

Cả hai họ mô hình đều khởi tạo từ trọng số tiền huấn luyện trên COCO \citep{lin2014coco} và được huấn luyện trên GPU Tesla T4 (Kaggle).

Các cấu hình được mô tả ở Mục~\ref{sec:method}. Phần này chỉ liệt kê siêu tham số.

\paragraph{YOLOv8s} (Ultralytics 8.4.204).
\begin{itemize}
\item 50 epoch, dừng sớm nếu 15 epoch liên tiếp không cải thiện, seed 0.
\item Batch 16 ở imgsz 640, batch 8 ở imgsz 1280.
\item Augmentation để mặc định: mosaic (tắt ở 10 epoch cuối), lật ngang với xác suất 0,5, co giãn $\pm 50\%$.
\item Cấu hình cắt ô: 40 epoch, batch 16. Một box chỉ được giữ trong ô nếu ít nhất 60\% diện tích nằm trong ô. Các ô không có biển chỉ giữ ngẫu nhiên 15\%.
\item Thời gian: mỗi lần huấn luyện ở imgsz 1280 mất khoảng 2,2 giờ trên T4.
\end{itemize}

\paragraph{Faster R-CNN} (torchvision).
\begin{itemize}
\item SGD: lr 0,0025, momentum 0,9, weight decay $5 \cdot 10^{-4}$.
\item 30 epoch, giảm lr 10 lần ở epoch 20 và 26.
\item Batch 8, FP16, cắt gradient ở chuẩn 1,0.
\end{itemize}
Thời gian huấn luyện của các cấu hình còn lại không được ghi lại.

\subsection{Chỉ số đánh giá}
\label{sec:metrics}

Mọi mô hình xuất dự đoán theo cùng một định dạng: tọa độ pixel trên ảnh gốc, ngưỡng độ tin cậy rất thấp (0,001), để không cắt mất phần cuối của đường precision--recall.
Tất cả được chấm bằng cùng một script dựa trên \texttt{pycocotools}.
Chỉ số chính là mAP@0.5. Nhóm báo cáo thêm mAP@0.5:0.95 và AP theo kích thước theo chuẩn COCO: nhỏ ($< 32^2$), vừa ($32^2$--$96^2$) và lớn ($> 96^2$ pixel).
Tập val chỉ có 23 biển lớn, nên AP\textsubscript{L} trên val dao động mạnh và không được dùng để so sánh.
Để biết chênh lệch giữa hai mô hình có ý nghĩa thống kê hay không, nhóm dùng \emph{paired bootstrap} \citep{efron1993bootstrap}: lấy mẫu lại có hoàn lại các ảnh của tập đánh giá 1000 lần, mỗi lần tính lại AP50 của các mô hình trên cùng bộ ảnh, rồi lấy khoảng 2,5\%--97,5\% của chênh lệch làm khoảng tin cậy 95\%.

\subsection{Kết quả}
\label{sec:results}

\begin{table*}[t]
\centering
\small
\begin{tabular}{llcccccccc}
\toprule
& & \multicolumn{4}{c}{Val (696 ảnh)} & \multicolumn{4}{c}{Test (871 ảnh)} \\
\cmidrule(lr){3-6} \cmidrule(lr){7-10}
Mô hình & Cấu hình & mAP@.5 & mAP@.5:.95 & AP\textsubscript{S} & AP\textsubscript{M} & mAP@.5 & mAP@.5:.95 & AP\textsubscript{S} & AP\textsubscript{M} \\
\midrule
\multirow{4}{*}{YOLOv8s}
& 640 (baseline)    & 61,7 & 39,8 & 29,8 & 70,5 & -- & -- & -- & -- \\
& 640 + SAHI        & 78,1 & 45,2 & 37,3 & 70,8 & -- & -- & -- & -- \\
& 1280              & 82,2 & 56,6 & 50,2 & \textbf{76,1} & 76,8 & 50,8 & 45,1 & \textbf{76,9} \\
& Cắt ô 640 + SAHI  & \textbf{83,9} & \textbf{57,9} & \textbf{52,3} & 75,2 & \textbf{78,1} & \textbf{51,4} & \textbf{46,8} & 75,1 \\
\midrule
\multirow{3}{*}{\shortstack[l]{Faster R-CNN\\R50-FPN}}
& exp01 (baseline)   & 64,8 & 41,6 & 33,5 & 68,1 & -- & -- & -- & -- \\
& exp02 (anchor nhỏ) & 72,5 & 47,7 & 40,3 & 71,2 & -- & -- & -- & -- \\
& exp03 (+ augment)  & 75,2 & 48,6 & 41,3 & 71,8 & 69,3 & 41,8 & 35,9 & 71,4 \\
\bottomrule
\end{tabular}
\caption{Kết quả trên val và test (\%). Chỉ các mô hình cuối cùng được chấm trên test, và chỉ chấm một lần.
Dự đoán val của Faster R-CNN được xuất ở ngưỡng 0,05 thay vì 0,001, nên mAP trên val của Faster R-CNN có thể thấp hơn thực tế một chút; dự đoán test được xuất ở ngưỡng 0,0001.}
\label{tab:results}
\end{table*}

Bảng~\ref{tab:results} tổng hợp kết quả. Ba nhận xét chính:

\paragraph{Biển nhỏ quyết định kết quả.}
Ở cấu hình baseline, AP của biển nhỏ kém biển vừa 40,7 điểm với YOLO (ảnh 640) và 34,6 điểm với Faster R-CNN.
Các cách giữ được độ phân giải chủ yếu cải thiện biển nhỏ. Với YOLO, chuyển từ 640 lên 1280 tăng AP\textsubscript{S} từ 29,8 lên 50,2 (+20,4), trong khi AP\textsubscript{M} chỉ tăng 5,6.
Với Faster R-CNN, anchor nhỏ và không thu nhỏ ảnh (exp02) tăng AP\textsubscript{S} thêm 6,8 điểm.
SAHI áp lên mô hình 640 có sẵn (không huấn luyện lại) đã tăng mAP@0.5 thêm 16,4 điểm. Nhưng mAP@0.5:0.95 chỉ tăng 5,4 điểm, nghĩa là box ghép từ các ô kém khít hơn.

\paragraph{YOLOv8s tốt hơn Faster R-CNN rõ rệt.}
Trên test, YOLOv8s 1280 hơn Faster R-CNN exp03 7,5 điểm mAP@0.5. Khoảng tin cậy 95\% của chênh lệch là $[4{,}5; 10{,}2]$, tức có ý nghĩa thống kê.
Faster R-CNN kém hơn ở cả 7 lớp, trong đó 4 lớp có chênh lệch có ý nghĩa thống kê. Kém rõ nhất là \textit{Cấm còn lại} ($-20{,}2$) và \textit{Cấm rẽ} ($-10{,}2$).

\paragraph{Cắt ô + SAHI và 1280 ngang nhau.}
Trên test, cắt ô + SAHI hơn 1280 1,3 điểm mAP@0.5. Nhưng khoảng tin cậy $[-0{,}5; +3{,}2]$ chứa 0, nên chưa đủ bằng chứng để nói cách nào tốt hơn.
Cắt ô + SAHI tốt hơn ở \textit{Hiệu lệnh} (+7,7) và \textit{Cấm còn lại} (+5,1), nhưng tạo nhiều box lặp hơn (78 so với 17) và báo nhầm nền nhiều hơn (123 so với 97 box).
Vì vậy cả hai cấu hình đều được giữ làm kết quả cuối. Cấu hình 1280 đơn giản hơn và suy luận nhanh hơn.

\paragraph{Từ val sang test.}
mAP@0.5 của cả ba mô hình cuối đều giảm 5--6 điểm từ val sang test, và thứ hạng giữa chúng không đổi.
Tập test khó hơn val: 21\% biển trong test nhỏ hơn 8 pixel, trong khi ở val chỉ có 10\%. Ở cả ba mô hình, nhóm biển dưới 8 pixel bị bỏ sót 60--78\%.
Chọn cấu hình trên val vẫn dẫn tới cùng kết luận trên test.

\begin{table}[t]
\centering
\small
\begin{tabular}{lrccc}
\toprule
Lớp & Biển & 1280 & Cắt ô+SAHI & FRCNN \\
\midrule
Cấm ngược chiều & 253 & 84,7 & \textbf{86,8} & 84,2 \\
Cấm dừng và đỗ  & 209 & 87,9 & \textbf{90,1} & 82,7 \\
Cấm rẽ          &  81 & \textbf{64,6} & 58,9 & 54,4 \\
Giới hạn tốc độ &  79 & \textbf{65,4} & 61,7 & 58,4 \\
Cấm còn lại     & 176 & 70,0 & \textbf{75,1} & 49,9 \\
Nguy hiểm       & 690 & 87,6 & \textbf{89,1} & 81,1 \\
Hiệu lệnh       & 260 & 77,5 & \textbf{85,3} & 74,7 \\
\bottomrule
\end{tabular}
\caption{AP50 (\%) theo lớp trên test. FRCNN là Faster R-CNN exp03.}
\label{tab:per-class-test}
\end{table}

Theo Bảng~\ref{tab:per-class-test}, ba lớp biển tròn viền đỏ (\textit{Cấm rẽ}, \textit{Giới hạn tốc độ}, \textit{Cấm còn lại}) có AP thấp nhất ở cả ba mô hình. \textit{Cấm ngược chiều} thì ngược lại: lớp này có tới 90,5\% box nhỏ nhưng vẫn đạt AP cao. Mục~
ef{sec:confusion} giải thích hai hiện tượng này.

% Cần: booktabs, multirow (\usepackage{multirow}), tham khảo trong custom.bib (zalo2020traffic, lin2014coco, jocher2023ultralytics,
% ren2015faster, lin2017fpn, akyon2022sahi, efron1993bootstrap)
% ------------------------------------------------

\section{Thực nghiệm}
\label{sec:experiments}

\subsection{Dữ liệu và cách chia tập}
\label{sec:data-split}

\paragraph{Gộp box trùng.}
Khi kiểm tra nhãn bằng mắt, nhóm phát hiện nhiều biển bị gán 2--3 box gần như trùng nhau, cùng lớp, chỉ lệch 1--2 pixel. Nếu giữ nguyên, mỗi biển chỉ khớp được với một dự đoán, nên các box lặp luôn bị tính là bỏ sót. Với tập test ban đầu, recall tối đa chỉ còn khoảng 77\% kể cả khi mô hình hoàn hảo.
Vì vậy, trong mỗi ảnh, nhóm gộp các box \emph{cùng lớp} có IoU $\geq 0{,}45$ thành một box có tọa độ trung bình. Các box được nối bắc cầu: A trùng B, B trùng C thì A, B, C chung một cụm.
Ngưỡng này an toàn vì hai biển thật đứng cạnh nhau gần như không chồng lên nhau (IoU $\approx 0$): 124/125 cặp có IoU trong khoảng $[0{,}45; 0{,}7)$ là box lặp cùng lớp.
Sau khi gộp, số box giảm từ 11.000 xuống 9.077 (17,5\%). Tập test bị nặng nhất: 24,7\% box bị bỏ.

\paragraph{Đặc điểm dữ liệu.}
Sau khi gộp, 77,6\% box là box nhỏ theo chuẩn COCO (diện tích $< 32^2$ pixel). Cạnh trung vị $\sqrt{wh}$ chỉ 16,5 pixel trên ảnh rộng 1622 pixel.
Khi thu nhỏ ảnh về cạnh dài 640, cạnh trung vị còn 6,5 pixel, và 68,9\% box có cạnh dưới 10 pixel.
Lớp nhiều nhất (\textit{Nguy hiểm}, 2.329 box) nhiều gấp 4,8 lần lớp ít nhất (\textit{Cấm rẽ}, 489 box).
Mỗi lớp đã gộp cả biến thể trái và phải, nên lật ngang ảnh không làm đổi lớp.

\paragraph{Chia tập theo tuyến đường.}
Ảnh là các khung hình liên tiếp trên 38 tuyến đường (trường \texttt{street\_id}). Nếu chia ngẫu nhiên theo ảnh, hai khung hình gần như giống hệt nhau có thể rơi vào cả train lẫn tập đánh giá.
Vì vậy nhóm chia theo tuyến đường, tỉ lệ 70/15/15 (seed 42): mọi ảnh của một tuyến đường nằm trọn trong một tập. Các tuyến đường được phân tầng theo lớp hiếm nhất có trong tuyến, để lớp hiếm có mặt ở cả ba tập (Bảng~\ref{tab:split}).
Mục~
ef{sec:leakage} định lượng mức đánh giá quá cao nếu chia ngẫu nhiên.
Mọi mô hình dùng chung bộ chia này. Chọn cấu hình chỉ dựa trên tập val. Tập test chỉ được chấm \emph{một lần}, cho các mô hình cuối cùng.

\begin{table}[t]
\centering
\small
\begin{tabular}{lrrr}
\toprule
Tập & Ảnh & Box & Tuyến đường \\
\midrule
Train & 2.933 & 5.693 & 28 \\
Val   &   696 & 1.636 &  5 \\
Test  &   871 & 1.748 &  5 \\
\bottomrule
\end{tabular}
\caption{Ba tập dữ liệu sau khi gộp box trùng, chia theo tuyến đường.}
\label{tab:split}
\end{table}

\subsection{Thiết lập huấn luyện}
\label{sec:setup}

Cả hai họ mô hình đều khởi tạo từ trọng số tiền huấn luyện trên COCO \citep{lin2014coco} và được huấn luyện trên GPU Tesla T4 (Kaggle).

Các cấu hình được mô tả ở Mục~\ref{sec:method}. Phần này chỉ liệt kê siêu tham số.

\paragraph{YOLOv8s} (Ultralytics 8.4.204).
\begin{itemize}
\item 50 epoch, dừng sớm nếu 15 epoch liên tiếp không cải thiện, seed 0.
\item Batch 16 ở imgsz 640, batch 8 ở imgsz 1280.
\item Augmentation để mặc định: mosaic (tắt ở 10 epoch cuối), lật ngang với xác suất 0,5, co giãn $\pm 50\%$.
\item Cấu hình cắt ô: 40 epoch, batch 16. Một box chỉ được giữ trong ô nếu ít nhất 60\% diện tích nằm trong ô. Các ô không có biển chỉ giữ ngẫu nhiên 15\%.
\item Thời gian: mỗi lần huấn luyện ở imgsz 1280 mất khoảng 2,2 giờ trên T4.
\end{itemize}

\paragraph{Faster R-CNN} (torchvision).
\begin{itemize}
\item SGD: lr 0,0025, momentum 0,9, weight decay $5 \cdot 10^{-4}$.
\item 30 epoch, giảm lr 10 lần ở epoch 20 và 26.
\item Batch 8, FP16, cắt gradient ở chuẩn 1,0.
\end{itemize}
Thời gian huấn luyện của các cấu hình còn lại không được ghi lại.

\subsection{Chỉ số đánh giá}
\label{sec:metrics}

Mọi mô hình xuất dự đoán theo cùng một định dạng: tọa độ pixel trên ảnh gốc, ngưỡng độ tin cậy rất thấp (0,001), để không cắt mất phần cuối của đường precision--recall.
Tất cả được chấm bằng cùng một script dựa trên \texttt{pycocotools}.
Chỉ số chính là mAP@0.5. Nhóm báo cáo thêm mAP@0.5:0.95 và AP theo kích thước theo chuẩn COCO: nhỏ ($< 32^2$), vừa ($32^2$--$96^2$) và lớn ($> 96^2$ pixel).
Tập val chỉ có 23 biển lớn, nên AP\textsubscript{L} trên val dao động mạnh và không được dùng để so sánh.
Để biết chênh lệch giữa hai mô hình có ý nghĩa thống kê hay không, nhóm dùng \emph{paired bootstrap} \citep{efron1993bootstrap}: lấy mẫu lại có hoàn lại các ảnh của tập đánh giá 1000 lần, mỗi lần tính lại AP50 của các mô hình trên cùng bộ ảnh, rồi lấy khoảng 2,5\%--97,5\% của chênh lệch làm khoảng tin cậy 95\%.

\subsection{Kết quả}
\label{sec:results}

\begin{table*}[t]
\centering
\small
\begin{tabular}{llcccccccc}
\toprule
& & \multicolumn{4}{c}{Val (696 ảnh)} & \multicolumn{4}{c}{Test (871 ảnh)} \\
\cmidrule(lr){3-6} \cmidrule(lr){7-10}
Mô hình & Cấu hình & mAP@.5 & mAP@.5:.95 & AP\textsubscript{S} & AP\textsubscript{M} & mAP@.5 & mAP@.5:.95 & AP\textsubscript{S} & AP\textsubscript{M} \\
\midrule
\multirow{4}{*}{YOLOv8s}
& 640 (baseline)    & 61,7 & 39,8 & 29,8 & 70,5 & -- & -- & -- & -- \\
& 640 + SAHI        & 78,1 & 45,2 & 37,3 & 70,8 & -- & -- & -- & -- \\
& 1280              & 82,2 & 56,6 & 50,2 & \textbf{76,1} & 76,8 & 50,8 & 45,1 & \textbf{76,9} \\
& Cắt ô 640 + SAHI  & \textbf{83,9} & \textbf{57,9} & \textbf{52,3} & 75,2 & \textbf{78,1} & \textbf{51,4} & \textbf{46,8} & 75,1 \\
\midrule
\multirow{3}{*}{\shortstack[l]{Faster R-CNN\\R50-FPN}}
& exp01 (baseline)   & 64,8 & 41,6 & 33,5 & 68,1 & -- & -- & -- & -- \\
& exp02 (anchor nhỏ) & 72,5 & 47,7 & 40,3 & 71,2 & -- & -- & -- & -- \\
& exp03 (+ augment)  & 75,2 & 48,6 & 41,3 & 71,8 & 69,3 & 41,8 & 35,9 & 71,4 \\
\bottomrule
\end{tabular}
\caption{Kết quả trên val và test (\%). Chỉ các mô hình cuối cùng được chấm trên test, và chỉ chấm một lần.
Dự đoán val của Faster R-CNN được xuất ở ngưỡng 0,05 thay vì 0,001, nên mAP trên val của Faster R-CNN có thể thấp hơn thực tế một chút; dự đoán test được xuất ở ngưỡng 0,0001.}
\label{tab:results}
\end{table*}

Bảng~\ref{tab:results} tổng hợp kết quả. Ba nhận xét chính:

\paragraph{Biển nhỏ quyết định kết quả.}
Ở cấu hình baseline, AP của biển nhỏ kém biển vừa 40,7 điểm với YOLO (ảnh 640) và 34,6 điểm với Faster R-CNN.
Các cách giữ được độ phân giải chủ yếu cải thiện biển nhỏ. Với YOLO, chuyển từ 640 lên 1280 tăng AP\textsubscript{S} từ 29,8 lên 50,2 (+20,4), trong khi AP\textsubscript{M} chỉ tăng 5,6.
Với Faster R-CNN, anchor nhỏ và không thu nhỏ ảnh (exp02) tăng AP\textsubscript{S} thêm 6,8 điểm.
SAHI áp lên mô hình 640 có sẵn (không huấn luyện lại) đã tăng mAP@0.5 thêm 16,4 điểm. Nhưng mAP@0.5:0.95 chỉ tăng 5,4 điểm, nghĩa là box ghép từ các ô kém khít hơn.

\paragraph{YOLOv8s tốt hơn Faster R-CNN rõ rệt.}
Trên test, YOLOv8s 1280 hơn Faster R-CNN exp03 7,5 điểm mAP@0.5. Khoảng tin cậy 95\% của chênh lệch là $[4{,}5; 10{,}2]$, tức có ý nghĩa thống kê.
Faster R-CNN kém hơn ở cả 7 lớp, trong đó 4 lớp có chênh lệch có ý nghĩa thống kê. Kém rõ nhất là \textit{Cấm còn lại} ($-20{,}2$) và \textit{Cấm rẽ} ($-10{,}2$).

\paragraph{Cắt ô + SAHI và 1280 ngang nhau.}
Trên test, cắt ô + SAHI hơn 1280 1,3 điểm mAP@0.5. Nhưng khoảng tin cậy $[-0{,}5; +3{,}2]$ chứa 0, nên chưa đủ bằng chứng để nói cách nào tốt hơn.
Cắt ô + SAHI tốt hơn ở \textit{Hiệu lệnh} (+7,7) và \textit{Cấm còn lại} (+5,1), nhưng tạo nhiều box lặp hơn (78 so với 17) và báo nhầm nền nhiều hơn (123 so với 97 box).
Vì vậy cả hai cấu hình đều được giữ làm kết quả cuối. Cấu hình 1280 đơn giản hơn và suy luận nhanh hơn.

\paragraph{Từ val sang test.}
mAP@0.5 của cả ba mô hình cuối đều giảm 5--6 điểm từ val sang test, và thứ hạng giữa chúng không đổi.
Tập test khó hơn val: 21\% biển trong test nhỏ hơn 8 pixel, trong khi ở val chỉ có 10\%. Ở cả ba mô hình, nhóm biển dưới 8 pixel bị bỏ sót 60--78\%.
Chọn cấu hình trên val vẫn dẫn tới cùng kết luận trên test.

\begin{table}[t]
\centering
\small
\begin{tabular}{lrccc}
\toprule
Lớp & Biển & 1280 & Cắt ô+SAHI & FRCNN \\
\midrule
Cấm ngược chiều & 253 & 84,7 & \textbf{86,8} & 84,2 \\
Cấm dừng và đỗ  & 209 & 87,9 & \textbf{90,1} & 82,7 \\
Cấm rẽ          &  81 & \textbf{64,6} & 58,9 & 54,4 \\
Giới hạn tốc độ &  79 & \textbf{65,4} & 61,7 & 58,4 \\
Cấm còn lại     & 176 & 70,0 & \textbf{75,1} & 49,9 \\
Nguy hiểm       & 690 & 87,6 & \textbf{89,1} & 81,1 \\
Hiệu lệnh       & 260 & 77,5 & \textbf{85,3} & 74,7 \\
\bottomrule
\end{tabular}
\caption{AP50 (\%) theo lớp trên test. FRCNN là Faster R-CNN exp03.}
\label{tab:per-class-test}
\end{table}

Theo Bảng~\ref{tab:per-class-test}, ba lớp biển tròn viền đỏ (\textit{Cấm rẽ}, \textit{Giới hạn tốc độ}, \textit{Cấm còn lại}) có AP thấp nhất ở cả ba mô hình. \textit{Cấm ngược chiều} thì ngược lại: lớp này có tới 90,5\% box nhỏ nhưng vẫn đạt AP cao. Mục~
ef{sec:confusion} giải thích hai hiện tượng này.

% Cần: booktabs, multirow (\usepackage{multirow}), tham khảo trong custom.bib (zalo2020traffic, lin2014coco, jocher2023ultralytics,
% ren2015faster, lin2017fpn, akyon2022sahi, efron1993bootstrap)
% ------------------------------------------------

\section{Thực nghiệm}
\label{sec:experiments}

\subsection{Dữ liệu và cách chia tập}
\label{sec:data-split}

\paragraph{Gộp box trùng.}
Khi kiểm tra nhãn bằng mắt, nhóm phát hiện nhiều biển bị gán 2--3 box gần như trùng nhau, cùng lớp, chỉ lệch 1--2 pixel. Nếu giữ nguyên, mỗi biển chỉ khớp được với một dự đoán, nên các box lặp luôn bị tính là bỏ sót. Với tập test ban đầu, recall tối đa chỉ còn khoảng 77\% kể cả khi mô hình hoàn hảo.
Vì vậy, trong mỗi ảnh, nhóm gộp các box \emph{cùng lớp} có IoU $\geq 0{,}45$ thành một box có tọa độ trung bình. Các box được nối bắc cầu: A trùng B, B trùng C thì A, B, C chung một cụm.
Ngưỡng này an toàn vì hai biển thật đứng cạnh nhau gần như không chồng lên nhau (IoU $\approx 0$): 124/125 cặp có IoU trong khoảng $[0{,}45; 0{,}7)$ là box lặp cùng lớp.
Sau khi gộp, số box giảm từ 11.000 xuống 9.077 (17,5\%). Tập test bị nặng nhất: 24,7\% box bị bỏ.

\paragraph{Đặc điểm dữ liệu.}
Sau khi gộp, 77,6\% box là box nhỏ theo chuẩn COCO (diện tích $< 32^2$ pixel). Cạnh trung vị $\sqrt{wh}$ chỉ 16,5 pixel trên ảnh rộng 1622 pixel.
Khi thu nhỏ ảnh về cạnh dài 640, cạnh trung vị còn 6,5 pixel, và 68,9\% box có cạnh dưới 10 pixel.
Lớp nhiều nhất (\textit{Nguy hiểm}, 2.329 box) nhiều gấp 4,8 lần lớp ít nhất (\textit{Cấm rẽ}, 489 box).
Mỗi lớp đã gộp cả biến thể trái và phải, nên lật ngang ảnh không làm đổi lớp.

\paragraph{Chia tập theo tuyến đường.}
Ảnh là các khung hình liên tiếp trên 38 tuyến đường (trường \texttt{street\_id}). Nếu chia ngẫu nhiên theo ảnh, hai khung hình gần như giống hệt nhau có thể rơi vào cả train lẫn tập đánh giá.
Vì vậy nhóm chia theo tuyến đường, tỉ lệ 70/15/15 (seed 42): mọi ảnh của một tuyến đường nằm trọn trong một tập. Các tuyến đường được phân tầng theo lớp hiếm nhất có trong tuyến, để lớp hiếm có mặt ở cả ba tập (Bảng~\ref{tab:split}).
Mục~
ef{sec:leakage} định lượng mức đánh giá quá cao nếu chia ngẫu nhiên.
Mọi mô hình dùng chung bộ chia này. Chọn cấu hình chỉ dựa trên tập val. Tập test chỉ được chấm \emph{một lần}, cho các mô hình cuối cùng.

\begin{table}[t]
\centering
\small
\begin{tabular}{lrrr}
\toprule
Tập & Ảnh & Box & Tuyến đường \\
\midrule
Train & 2.933 & 5.693 & 28 \\
Val   &   696 & 1.636 &  5 \\
Test  &   871 & 1.748 &  5 \\
\bottomrule
\end{tabular}
\caption{Ba tập dữ liệu sau khi gộp box trùng, chia theo tuyến đường.}
\label{tab:split}
\end{table}

\subsection{Thiết lập huấn luyện}
\label{sec:setup}

Cả hai họ mô hình đều khởi tạo từ trọng số tiền huấn luyện trên COCO \citep{lin2014coco} và được huấn luyện trên GPU Tesla T4 (Kaggle).

Các cấu hình được mô tả ở Mục~\ref{sec:method}. Phần này chỉ liệt kê siêu tham số.

\paragraph{YOLOv8s} (Ultralytics 8.4.204).
\begin{itemize}
\item 50 epoch, dừng sớm nếu 15 epoch liên tiếp không cải thiện, seed 0.
\item Batch 16 ở imgsz 640, batch 8 ở imgsz 1280.
\item Augmentation để mặc định: mosaic (tắt ở 10 epoch cuối), lật ngang với xác suất 0,5, co giãn $\pm 50\%$.
\item Cấu hình cắt ô: 40 epoch, batch 16. Một box chỉ được giữ trong ô nếu ít nhất 60\% diện tích nằm trong ô. Các ô không có biển chỉ giữ ngẫu nhiên 15\%.
\item Thời gian: mỗi lần huấn luyện ở imgsz 1280 mất khoảng 2,2 giờ trên T4.
\end{itemize}

\paragraph{Faster R-CNN} (torchvision).
\begin{itemize}
\item SGD: lr 0,0025, momentum 0,9, weight decay $5 \cdot 10^{-4}$.
\item 30 epoch, giảm lr 10 lần ở epoch 20 và 26.
\item Batch 8, FP16, cắt gradient ở chuẩn 1,0.
\end{itemize}
Thời gian huấn luyện của các cấu hình còn lại không được ghi lại.

\subsection{Chỉ số đánh giá}
\label{sec:metrics}

Mọi mô hình xuất dự đoán theo cùng một định dạng: tọa độ pixel trên ảnh gốc, ngưỡng độ tin cậy rất thấp (0,001), để không cắt mất phần cuối của đường precision--recall.
Tất cả được chấm bằng cùng một script dựa trên \texttt{pycocotools}.
Chỉ số chính là mAP@0.5. Nhóm báo cáo thêm mAP@0.5:0.95 và AP theo kích thước theo chuẩn COCO: nhỏ ($< 32^2$), vừa ($32^2$--$96^2$) và lớn ($> 96^2$ pixel).
Tập val chỉ có 23 biển lớn, nên AP\textsubscript{L} trên val dao động mạnh và không được dùng để so sánh.
Để biết chênh lệch giữa hai mô hình có ý nghĩa thống kê hay không, nhóm dùng \emph{paired bootstrap} \citep{efron1993bootstrap}: lấy mẫu lại có hoàn lại các ảnh của tập đánh giá 1000 lần, mỗi lần tính lại AP50 của các mô hình trên cùng bộ ảnh, rồi lấy khoảng 2,5\%--97,5\% của chênh lệch làm khoảng tin cậy 95\%.

\subsection{Kết quả}
\label{sec:results}

\begin{table*}[t]
\centering
\small
\begin{tabular}{llcccccccc}
\toprule
& & \multicolumn{4}{c}{Val (696 ảnh)} & \multicolumn{4}{c}{Test (871 ảnh)} \\
\cmidrule(lr){3-6} \cmidrule(lr){7-10}
Mô hình & Cấu hình & mAP@.5 & mAP@.5:.95 & AP\textsubscript{S} & AP\textsubscript{M} & mAP@.5 & mAP@.5:.95 & AP\textsubscript{S} & AP\textsubscript{M} \\
\midrule
\multirow{4}{*}{YOLOv8s}
& 640 (baseline)    & 61,7 & 39,8 & 29,8 & 70,5 & -- & -- & -- & -- \\
& 640 + SAHI        & 78,1 & 45,2 & 37,3 & 70,8 & -- & -- & -- & -- \\
& 1280              & 82,2 & 56,6 & 50,2 & \textbf{76,1} & 76,8 & 50,8 & 45,1 & \textbf{76,9} \\
& Cắt ô 640 + SAHI  & \textbf{83,9} & \textbf{57,9} & \textbf{52,3} & 75,2 & \textbf{78,1} & \textbf{51,4} & \textbf{46,8} & 75,1 \\
\midrule
\multirow{3}{*}{\shortstack[l]{Faster R-CNN\\R50-FPN}}
& exp01 (baseline)   & 64,8 & 41,6 & 33,5 & 68,1 & -- & -- & -- & -- \\
& exp02 (anchor nhỏ) & 72,5 & 47,7 & 40,3 & 71,2 & -- & -- & -- & -- \\
& exp03 (+ augment)  & 75,2 & 48,6 & 41,3 & 71,8 & 69,3 & 41,8 & 35,9 & 71,4 \\
\bottomrule
\end{tabular}
\caption{Kết quả trên val và test (\%). Chỉ các mô hình cuối cùng được chấm trên test, và chỉ chấm một lần.
Dự đoán val của Faster R-CNN được xuất ở ngưỡng 0,05 thay vì 0,001, nên mAP trên val của Faster R-CNN có thể thấp hơn thực tế một chút; dự đoán test được xuất ở ngưỡng 0,0001.}
\label{tab:results}
\end{table*}

Bảng~\ref{tab:results} tổng hợp kết quả. Ba nhận xét chính:

\paragraph{Biển nhỏ quyết định kết quả.}
Ở cấu hình baseline, AP của biển nhỏ kém biển vừa 40,7 điểm với YOLO (ảnh 640) và 34,6 điểm với Faster R-CNN.
Các cách giữ được độ phân giải chủ yếu cải thiện biển nhỏ. Với YOLO, chuyển từ 640 lên 1280 tăng AP\textsubscript{S} từ 29,8 lên 50,2 (+20,4), trong khi AP\textsubscript{M} chỉ tăng 5,6.
Với Faster R-CNN, anchor nhỏ và không thu nhỏ ảnh (exp02) tăng AP\textsubscript{S} thêm 6,8 điểm.
SAHI áp lên mô hình 640 có sẵn (không huấn luyện lại) đã tăng mAP@0.5 thêm 16,4 điểm. Nhưng mAP@0.5:0.95 chỉ tăng 5,4 điểm, nghĩa là box ghép từ các ô kém khít hơn.

\paragraph{YOLOv8s tốt hơn Faster R-CNN rõ rệt.}
Trên test, YOLOv8s 1280 hơn Faster R-CNN exp03 7,5 điểm mAP@0.5. Khoảng tin cậy 95\% của chênh lệch là $[4{,}5; 10{,}2]$, tức có ý nghĩa thống kê.
Faster R-CNN kém hơn ở cả 7 lớp, trong đó 4 lớp có chênh lệch có ý nghĩa thống kê. Kém rõ nhất là \textit{Cấm còn lại} ($-20{,}2$) và \textit{Cấm rẽ} ($-10{,}2$).

\paragraph{Cắt ô + SAHI và 1280 ngang nhau.}
Trên test, cắt ô + SAHI hơn 1280 1,3 điểm mAP@0.5. Nhưng khoảng tin cậy $[-0{,}5; +3{,}2]$ chứa 0, nên chưa đủ bằng chứng để nói cách nào tốt hơn.
Cắt ô + SAHI tốt hơn ở \textit{Hiệu lệnh} (+7,7) và \textit{Cấm còn lại} (+5,1), nhưng tạo nhiều box lặp hơn (78 so với 17) và báo nhầm nền nhiều hơn (123 so với 97 box).
Vì vậy cả hai cấu hình đều được giữ làm kết quả cuối. Cấu hình 1280 đơn giản hơn và suy luận nhanh hơn.

\paragraph{Từ val sang test.}
mAP@0.5 của cả ba mô hình cuối đều giảm 5--6 điểm từ val sang test, và thứ hạng giữa chúng không đổi.
Tập test khó hơn val: 21\% biển trong test nhỏ hơn 8 pixel, trong khi ở val chỉ có 10\%. Ở cả ba mô hình, nhóm biển dưới 8 pixel bị bỏ sót 60--78\%.
Chọn cấu hình trên val vẫn dẫn tới cùng kết luận trên test.

\begin{table}[t]
\centering
\small
\begin{tabular}{lrccc}
\toprule
Lớp & Biển & 1280 & Cắt ô+SAHI & FRCNN \\
\midrule
Cấm ngược chiều & 253 & 84,7 & \textbf{86,8} & 84,2 \\
Cấm dừng và đỗ  & 209 & 87,9 & \textbf{90,1} & 82,7 \\
Cấm rẽ          &  81 & \textbf{64,6} & 58,9 & 54,4 \\
Giới hạn tốc độ &  79 & \textbf{65,4} & 61,7 & 58,4 \\
Cấm còn lại     & 176 & 70,0 & \textbf{75,1} & 49,9 \\
Nguy hiểm       & 690 & 87,6 & \textbf{89,1} & 81,1 \\
Hiệu lệnh       & 260 & 77,5 & \textbf{85,3} & 74,7 \\
\bottomrule
\end{tabular}
\caption{AP50 (\%) theo lớp trên test. FRCNN là Faster R-CNN exp03.}
\label{tab:per-class-test}
\end{table}

Theo Bảng~\ref{tab:per-class-test}, ba lớp biển tròn viền đỏ (\textit{Cấm rẽ}, \textit{Giới hạn tốc độ}, \textit{Cấm còn lại}) có AP thấp nhất ở cả ba mô hình. \textit{Cấm ngược chiều} thì ngược lại: lớp này có tới 90,5\% box nhỏ nhưng vẫn đạt AP cao. Mục~
ef{sec:confusion} giải thích hai hiện tượng này.
